% =====================================================================================================
% Introduction (no heading, AER style). Writer: section writer "introduction", 2026-09-24.
% Numbers verified against output/tables/*.csv; sources listed in paper/notes/writer_introduction.md.
% =====================================================================================================

Training a frontier language model has become one of the largest single investments in the technology sector.
The compute used to train frontier models has grown by a factor of four to five per year
\citep{sevilla2024training}, and the amortized hardware and energy cost of the largest training runs by a factor
of about 2.4 per year since 2016 \citep{cottier2024rising}. The planning tool for this investment is the neural
scaling law: an empirical relation between a model's loss $L$ on held-out text and the number of parameters $N$,
the number of training tokens $D$, and training compute $C \approx 6ND$ floating-point operations (FLOP)
\citep{hestness2017deep, kaplan2020scaling, hoffmann2022training}. The workhorse specification,
$L = E + AN^{-\alpha} + BD^{-\beta}$ with irreducible loss $E$ and positive constants $A$, $B$, $\alpha$ and
$\beta$ \citep{hoffmann2022training}, is used to decide how large a model to train, on how many tokens, and what
the next order of magnitude of compute will buy.

This paper starts from the observation that a scaling law is a production function and that fitting one is
production-function estimation. Loss is an inverse output index, parameters and tokens are inputs, iso-loss curves
are isoquants, and the ``IsoFLOP'' curves along which practitioners vary model size at fixed compute are isocost
lines. Economists have studied the estimation of production functions for eighty years
\citep{marschak1944random, griliches1998production, ackerberg2007econometric}, and many of the difficulties that
the machine-learning literature has met in fitting scaling laws are instances of problems that literature has named
and, in part, solved. We develop the mapping systematically and use it to obtain three sets of results.

\textit{What scaling data identify.}---Because training compute is multiplicative in the two inputs, cost
minimization equates the output elasticities of parameters and data, $\varepsilon_N = \varepsilon_D$, and no
price data are needed to exploit the first-order condition. The three estimation approaches of
\citet{hoffmann2022training} are, respectively, a cost-function estimator, a conditional-factor-demand estimator and
a primal estimator, linked by the cross-equation restrictions of duality (Proposition~\ref{prop:duality}). The
central identification problem is the functional dependence that \citet{ackerberg2015identification} diagnosed in
firm data, generated here by behavior rather than by the experimenter's design. A lab that minimizes training
compute places every run on its expansion path, where parameters and tokens are deterministic functions of compute
and the log of reducible loss is exactly linear, so that all information about curvature lies in deviations from
the path. Such data from a single lab identify the expansion path and the
loss--compute frontier, including the frontier elasticity $\gamma = \alpha\beta/(\alpha+\beta)$. They identify the curvature of the
technology only through functional form, and the compute-optimal ratio of tokens to parameters, $M^*$, only by
assuming that the observed path is optimal (Proposition~\ref{prop:fdep}). We show that the information about $M^*$
is of second order, and the information about curvature of fourth order, in the size of runs' log deviations from
the path (Proposition~\ref{prop:info}): the better labs optimize, the less their data reveal. Two further
results give economic content to the theory. Every interior compute optimum requires the elasticity of
substitution $\sigma$ between parameters and data to lie between zero and one (Lemma~\ref{lemma:sigma}), so the
existence of compute-optimal training is itself evidence of gross complementarity. And while the classic
non-identification of biased technical change \citep{diamond1978measurement} survives for its magnitude and for
$\sigma$, the \textit{sign} of the bias is identified by the drift of the compute-optimal tokens-per-parameter ratio at
fixed compute, measured on training-optimal allocations (Proposition~\ref{prop:dmr}).

These results have quantitative bite. In a Monte Carlo calibrated to the canonical Chinchilla data, with 90 runs
and equal total compute in every design, the root-mean-squared error of the standard estimator of the on-path elasticity of
substitution $\sigma^* = 2/(2+\alpha+\beta)$ is 0.28 when every run is compute-optimal and 0.004 in an IsoFLOP design. Once the
outer exponent $\kappa$ of the loss function is freed (the Chinchilla form imposes $\kappa=1$), the 95 percent
profile-likelihood interval for $\sigma^*$ covers the entire range $[0.50, 0.95]$ in 88 to 97 percent of
replications whenever allocation errors are small (a standard deviation of log tokens per parameter of 0.2 or less). The same pattern appears in the Chinchilla runs themselves: runs within 0.15 log
points of the fitted path yield a profile likelihood-ratio statistic of at most 2.0 over that range, whereas random
subsamples of the same size reach 79 to 101 and the full design reaches 424.

\textit{The technology.}---In each of seven public training sweeps, parameters and data are gross
complements.\footnote{The sweeps are the reconstructed Chinchilla runs \citep{hoffmann2022training,
besiroglu2024chinchilla}, Farseer \citep{li2025predictableb}, the three pretraining corpora of
\citet{gadre2024language}, the OLMo ladder \citep{bhagia2024establishing} and the single-epoch runs of
\citet{muennighoff2023scaling}; 25 DataDecide recipes \citep{magnusson2025datadecide} are used for the neutrality
tests. Section~\ref{sec:data} describes them.} Under the Chinchilla functional form, $\sigma^*$ lies between 0.73
and 0.83, with bootstrap standard errors between 0.006 and 0.034. But the form's restriction $\kappa = 1$ is rejected
in all seven sweeps, and when the data choose the curvature $\sigma^*$ falls to between 0.51 and 0.71. On Farseer's
near-factorial design of 404 runs, three estimators that do not impose the restriction agree on
$\sigma \approx 0.69$--$0.71$; the nonparametric median is 0.690, with a 95 percent confidence interval of
$[0.679, 0.699]$. The estimates that let the data choose the curvature overlap the top of the range typically
estimated for capital and labor
\citep{chirinko2008sigma, oberfield2021micro}. By contrast, the allocation exponent $a = \beta/(\alpha+\beta)$, which
governs how the compute-optimal model size $N^* \propto C^{a}$ scales, ranges from 0.37 to 0.57 across sweeps, and the
compute-optimal ratio $M^*$ at $10^{21}$ FLOP from 3 to 60 tokens per parameter. The object every allocation decision
needs is the one the data pin down least.

The lens also sorts two debates in the scaling literature. Better training data act mostly like a shift in the
asymptote plus a Hicks-neutral scaling of reducible loss: across the three corpora of \citet{gadre2024language} the
data do not reject that restriction ($p = 0.90$), and across the 25 DataDecide recipes asymptote shifts alone
explain 98 percent of the cross-recipe variation in fit. The remaining factor bias is systematic, however, and it
implies compute-optimal ratios that differ by a factor of 2.9--3.4 across recipes. The long-standing disagreement
between the allocation exponents of \citet{kaplan2020scaling} and \citet{hoffmann2022training}, which falls from 0.84
to 0.50 on the runs of \citet{porian2024resolving}, decomposes into input mismeasurement (39--47 percent), in the
sense of the spurious small-firm scale economies of \citet{nerlove1963returns}, and flexible inputs that were not
re-optimized at each scale (53--61 percent). In our own controlled experiment, two ``labs'' train the same model
sizes on corpora of different quality (FineWeb-Edu and FineWeb) under a shared tokenizer; \textcolor{red}{[TBD-m9:
one-sentence headline --- $\sigma$ by lab with $\kappa$ free, and whether data quality is Hicks-neutral or
factor-biased]}.

\textit{Revealed inference demand.}---The third result turns labs' choices into data. A model that will serve $T$
tokens of inference over its lifetime costs $6ND + 2NT$ FLOP to train and serve \citep{sardana2024beyond}. A lab that
minimizes this lifetime compute chooses inputs so that the ratio of output elasticities
$w \equiv \varepsilon_N/\varepsilon_D$ equals $1 + T/(3D)$, the ratio of lifetime to training compute
(Proposition~\ref{prop:wedge}). Inverting the first-order condition recovers planned inference demand from the degree
of over-training, much as \citet{deloecker2012markups} recover markups from the ratio of an output elasticity to a
cost share. The wedge here is not a markup, though: it is the shadow cost of a second use of the same input. Only
two features of the technology matter, $M^*$ and $\sigma^*$, because
$\ln w = (1/\sigma^* - 1)\ln\!\big(M/M^*(C)\big)$, where $M \equiv D/N$ and $M^*(C)$ is the training-optimal ratio
at the model's own compute. Among 173 open-weight base models released
between 2021 and 2025 with verified parameter and token counts, the median $w$ is 3.19: the median model is trained
as if its lifetime inference compute would be about 2.2 times its training compute. The sign is robust. Ninety-five
percent of models have $w > 1$, and for 79 percent the entire technology-sensitivity band, spanned by six
alternative estimates of the technology, lies above one. The level is not robust: the median $w$ ranges from 1.9 to
4.0 across those technologies.
Over-training is also recent. Among production-scale models (at least $10^{21}$ FLOP), the share with $w < 1$ fell from 47 percent in 2021 to
zero in 2024, and the median $w$ rose from 1.06 in 2022 to 4.68 in 2025. Llama 3 8B has $w = 5.27$ (95 percent
interval $[3.97, 8.12]$), about 13 tokens of planned lifetime inference per training token (6 under Meta's own
scaling law), while Meta's 405-billion-parameter flagship lies on Meta's own training-optimal path ($w = 0.98$). At
fixed compute and release date, log downloads rise about one for one with log tokens per parameter (0.97, standard
error 0.18; 0.72 with developer fixed effects). At fixed compute this is equivalent to smaller models being
downloaded more, so we read it as consistent with anticipated demand rather than as a measurement of $T$.

\textit{Observational scaling and ``algorithmic progress.''}---The same lens disciplines scaling claims drawn from
cross-lab data. In the spirit of \citet{lalonde1986evaluating}, we evaluate an experimentally estimated technology on
the observational design and compare estimators with the resulting benchmark. For 57 open models inside the
experimental support, pooled least squares recovers the experimental elasticity of benchmark (HellaSwag)
performance with respect to compute (bias $-0.001$, standard error 0.038, so biases of about 20 percent cannot be
ruled out), but the split between parameters and data is unreliable:
family fixed effects overstate the parameter elasticity by 0.20 and understate the data elasticity by 0.26. Much of
the apparent attenuation of data elasticities in cross-lab regressions is a design effect, because over-trained
models sit where the data elasticity is small; the standard remedies of the IO literature are not available in
public data, where candidate instruments are weak (first-stage $F \le 3.8$) or fail exclusion and there are too few
model generations for dynamic-panel estimators. We
replicate the estimate of \citet{ho2024algorithmic} exactly: compute requirements for a given loss halve every 8.4
months at the bootstrap median they report. The point estimate (8.7 months), however, is not the minimizer of their
objective (6.1 months at convergence, 95 percent interval $[3.0, 22.7]$); reasonable estimators of the same model
range from 6.1 to 10.2 months, and the profile-likelihood interval $[4.1, 40.5]$ traces a ridge of the kind
\citet{diamond1978measurement} predict. Finally, the rebalancing from the allocation rule of
\citet{kaplan2020scaling} to that of \citet{hoffmann2022training} is a gain in allocative, not technical, efficiency
\citep{farrell1957measurement}. Its realized gains among large models were 1.0--2.4-fold under the main
technologies, or 2--40 percent of the gain implied by the rate of \citet{ho2024algorithmic}; under technologies
estimated on other sweeps, even the sign of the gain changes. The tenfold gain reported by \citet{gundlach2025origin} is the counterfactual cost of
keeping Kaplan's rule at 2025 frontier compute.

\textit{Related literature.}---The paper connects three literatures. The first is the machine-learning literature
on neural scaling laws and their estimation. Power laws in model and data size \citep{hestness2017deep,
kaplan2020scaling} and the compute-optimal allocation of \citet{hoffmann2022training} have been refined for data
repetition \citep{muennighoff2023scaling}, over-training \citep{gadre2024language}, inference costs
\citep{sardana2024beyond}, hyperparameters \citep{li2025predictablea} and non-separable loss surfaces
\citep{li2025predictableb}. A parallel strand audits estimation practice: the replication of
\citet{besiroglu2024chinchilla}, surveys of fitting choices \citep{choshen2024hitchhikers, li2025misfitting}, the
finite-grid bias of IsoFLOP parabolas \citep{czech2026problems}, the reconciliations of the Kaplan and Chinchilla
exponents \citep{pearce2024reconciling, porian2024resolving}, and the ill-conditioning of designs with a fixed ratio
of tokens to parameters \citep{kricheli2026tokens}. We contribute a unifying framework in which these findings are
predictions. The collinearity of \citet{kricheli2026tokens} arises from design; ours arises from optimizing behavior,
and it explains why exponents and the frontier are precisely estimated in published fits while the levels $A$ and
$B$, and with them $M^*$, are not. The inference-aware
allocation problem of \citet{sardana2024beyond} is the forward problem whose inversion yields revealed demand. And
the Kaplan--Chinchilla reconciliation is input mismeasurement plus flexible inputs, with bias formulas that get the
sign and order of magnitude of each component right.

The second literature reads scaling through an economic lens. \citet{hao2026theory} derive the elasticity of
substitution implied by the Chinchilla form and study profit-maximizing training under a Leontief approximation. We
estimate $\sigma$ with standard errors across sweeps, derive its value on the compute-optimal path, and show that
$\sigma > 0$, which the Leontief approximation sets to zero, is what makes observed over-training optimal.
\citet{mertens2026secret} estimate developer-specific compute efficiency with fixed-effects regressions of benchmark
scores on compute; we add a design-matched experimental benchmark and document that model families trained on synthetic data top
the resulting productivity ranking, a gross-output problem. \citet{whitfill2025note} and
\citet{konig2026validity} warn that observational scaling estimates may be confounded; we supply the structure that
IO brings to that concern, in the form of timing assumptions, the sign of the bias under alternative behavioral
regimes, and bounds. Measurement of algorithmic progress \citep{erdil2022algorithmic, ho2024algorithmic,
gundlach2025origin} and observational scaling laws \citep{ruan2024observational, maiapolo2024sloth} are recast as
productivity measurement with the corresponding identification limits. Finally, models of AI and growth, automation
and LLM pricing calibrate the training technology rather than estimate it \citep{trammell2023economic,
erdil2025gate, bergemann2025economics}, and studies of market structure and of the inference market take it as
given \citep{korinek2025concentrating, demirer2025emerging}; our estimates of $\sigma$, $\gamma$ and $w$ are meant
as inputs to such models.

The third literature is the estimation of production functions in industrial organization: transmission bias
\citep{marschak1944random}, fixed effects \citep{mundlak1961empirical}, the timing and proxy approaches of
\citet{olley1996dynamics}, \citet{levinsohn2003estimating} and \citet{ackerberg2015identification}, first-order-condition
identification \citep{gandhi2020identification}, the production approach to wedges \citep{deloecker2012markups}
and its critiques \citep{bond2020unpleasant, raval2023testing}, factor-augmenting technical change
\citep{doraszelski2018measuring, demirer2020production}, the identification and normalization of the elasticity of
substitution \citep{diamond1978measurement, klump2007factor, leonledesma2010identifying}, cost-function duality and
measurement error in inputs \citep{nerlove1963returns, collardwexler2016production}, and efficiency and productivity
dispersion \citep{farrell1957measurement, foster2008reallocation, hsieh2009misallocation, syverson2011determines}.
The training of language models offers that literature an unusually favorable laboratory: the cost function is an
engineering identity, input prices drop out of the first-order conditions, and designed input variation exists at
scale in the form of IsoFLOP sweeps. This makes it possible to benchmark production-function estimators against
experiments in the manner of \citet{lalonde1986evaluating}, and to see which of the standard remedies survive.

We do not claim to be the first to read scaling laws as production functions, to compute an elasticity of
substitution between parameters and data, or to note that observational scaling estimates can be confounded. Our
contribution is the systematic mapping from scaling-law practice to the econometric toolkit of empirical IO, the
identification results that follow from it, estimates of the technology with inference across public sweeps, and the
inversion of labs' allocation choices into revealed inference demand.

The rest of the paper is organized as follows. Section~\ref{sec:framework} sets out the production problem, the
dictionary between scaling-law and IO concepts, and the duality and wedge results. Section~\ref{sec:ident} derives
what scaling data identify, with Monte Carlo evidence. Section~\ref{sec:data} describes the designed sweeps, our own
experiment and the observational data. Section~\ref{sec:tech} estimates the technology, and Section~\ref{sec:wedge}
recovers revealed inference demand. Section~\ref{sec:obs} turns to observational production functions and
algorithmic progress, and Section~\ref{sec:concl} concludes. Proofs are in Online Appendix A; data construction,
Monte Carlo designs and additional results are in Online Appendices B--D.
